\section{Validaciones}

Las validaciones se ejecutaron sobre Oracle Database 19c mediante DBeaver. Primero se validó la creación del schema, luego la creación de tablas, restricciones, vista, triggers y procedimiento. Posteriormente, se cargaron datos de prueba y se verificó el comportamiento del inventario y del registro de pagos.

\begin{figure}[H]
\centering
\includegraphics[width=\textwidth]{08_tablas_creadas_farmacias_bd3.png}
\caption{Tablas creadas en el schema FARMACIAS\_BD3.}
\label{fig:tablas-creadas}
\end{figure}

\begin{figure}[H]
\centering
\includegraphics[width=\textwidth]{09_triggers_procedimiento_validos.png}
\caption{Vista y triggers validados en Oracle.}
\label{fig:triggers-validos}
\end{figure}

\begin{figure}[H]
\centering
\includegraphics[width=\textwidth]{09a_procedimiento_registrar_pago_valido.png}
\caption{Procedimiento almacenado validado en Oracle.}
\label{fig:procedimiento-valido}
\end{figure}

\begin{figure}[H]
\centering
\includegraphics[width=\textwidth]{11_validacion_datos_prueba.png}
\caption{Validación de datos de prueba cargados.}
\label{fig:datos-prueba}
\end{figure}

\begin{figure}[H]
\centering
\includegraphics[width=\textwidth]{12_prueba_procedimiento_pago.png}
\caption{Prueba funcional del procedimiento de pago.}
\label{fig:prueba-pago}
\end{figure}
